\documentclass[11pt,a4paper]{article}
\usepackage[utf8]{inputenc}
\usepackage[T1]{fontenc}
\usepackage[spanish,es-nodecimaldot]{babel}
\usepackage{lmodern}
\usepackage[a4paper,margin=2.2cm]{geometry}
\usepackage{xcolor}
\usepackage{booktabs}
\usepackage{array}
\usepackage{longtable}
\usepackage{enumitem}
\usepackage{hyperref}
\usepackage{fancyhdr}
\definecolor{USMBlue}{HTML}{004B76}
\definecolor{USMLight}{HTML}{EAF3F8}
\definecolor{Muted}{HTML}{5C6B73}
\hypersetup{colorlinks=true,linkcolor=USMBlue,urlcolor=USMBlue,pdftitle={Narrativa de defensa de tesis},pdfauthor={Carlos Saldivia Heinz}}
\pagestyle{fancy}
\fancyhf{}
\lhead{Narrativa de defensa · Carlos Saldivia Heinz}
\rhead{\thepage}
\setlength{\headheight}{14pt}
\setlength{\parindent}{0pt}
\setlength{\parskip}{0.65em}
\hyphenpenalty=7000
\exhyphenpenalty=7000
\emergencystretch=2em
\newcommand{\Stage}[1]{\par\noindent\colorbox{USMLight}{\parbox{\dimexpr\linewidth-2\fboxsep\relax}{\small\textbf{Dirección escénica:} #1}}}
\newcommand{\NarrativeEntry}[7]{%
  \begin{minipage}{\linewidth}
  \raggedright
  \subsection*{Slide #1 · #2 \hfill \textcolor{USMBlue}{#3 · acumulado #4}}
  \addcontentsline{toc}{subsection}{Slide #1: #2 (#3)}
  #5
  \Stage{#6}
  {\small\color{Muted}\textbf{Respaldo:} #7\par}\medskip\hrule\medskip
  \end{minipage}\par
}
\newcommand{\QA}[4]{%
  \begin{minipage}{\linewidth}
  \raggedright
  \subsubsection*{#1}
  \noindent\textbf{Respuesta breve.} #2\par
  \noindent\textbf{Respuesta desarrollada.} #3\par
  {\small\color{Muted}\textbf{Respaldo:} #4\par}\medskip
  \end{minipage}\par
}
\title{\textbf{Narrativa para la defensa de título}\\[0.5em]\large Reconstrucción de Canales EEG Faltantes Mediante Procesamiento de Señales en Grafos con Regularización Temporal}
\author{Carlos Saldivia Heinz\\Universidad Técnica Federico Santa María\\Departamento de Electrónica}
\date{23 de julio de 2026}
\begin{document}
\maketitle
\begin{center}
\colorbox{USMLight}{\parbox{0.90\linewidth}{\textbf{Uso del documento.} El guion principal está calibrado a 29:00 minutos e incluye transiciones y pausas. Las respuestas del banco se ofrecen en versión breve y desarrollada; deben responderse primero en una frase y ampliarse solo si el jurado lo solicita.}}
\end{center}
\tableofcontents
\clearpage
\section*{Guion principal: 29 minutos}

\NarrativeEntry{1}{Portada y apertura}{0:30}{0:30}{
Buenos días. Esta defensa presenta un estudio sobre reconstrucción de canales EEG faltantes mediante procesamiento de señales en grafos y regularización temporal. La pregunta que organiza toda la exposición es simple de formular, pero exige evidencia cuidadosa: ¿cuándo la continuidad temporal agrega información útil a una referencia espacial madura como MNE-Python? Voy a responderla desde el problema, el diseño experimental, los resultados y, finalmente, las condiciones bajo las cuales conviene elegir cada método.
}{Mirar al jurado. Pausa breve después de la pregunta.}{Tesis final, portada e Introducción.}

\NarrativeEntry{2}{El problema real}{1:10}{1:40}{
Un canal EEG puede perderse por mal contacto, desplazamiento o artefactos persistentes. Descartar todo el registro reduce la información disponible; reconstruir el canal evita ese descarte, pero introduce una decisión metodológica antes de cualquier análisis posterior. Por eso no considero la interpolación como un simple relleno. Una reconstrucción modifica la topografía, la forma temporal y el contenido espectral que luego alimentan otras etapas. Para medir ese efecto con verdad de terreno, el estudio parte de señales completas y oculta canales de manera artificial. Así conocemos exactamente qué señal debió recuperarse y podemos comparar métodos bajo la misma máscara. Este control fortalece la validez interna, aunque no representa por completo un electrodo realmente dañado; esa frontera aparecerá nuevamente en las limitaciones.
}{Señalar la cadena de impacto de izquierda a derecha.}{Tesis, Introducción y protocolo de ocultamiento artificial.}

\NarrativeEntry{3}{Por qué no es trivial}{1:05}{2:45}{
La reconstrucción es difícil por cuatro razones. Primero, la información espacial depende de qué electrodos permanecen observados: una falla agrupada puede eliminar varios vecinos útiles al mismo tiempo. Segundo, el EEG tiene continuidad temporal, pero regularizarla en exceso puede suavizar componentes de interés. Tercero, las métricas no son equivalentes. Reducir el error absoluto no garantiza preservar mejor el espectro ni ejecutar más rápido. Cuarto, el método adecuado puede cambiar con la severidad y el patrón de pérdida. El problema, entonces, no es encontrar una técnica que gane siempre, sino caracterizar un régimen de decisión: cuándo la estructura espacio--temporal compensa su costo y cuándo una referencia espacial consolidada sigue siendo suficiente.
}{Marcar una pausa entre cada una de las cuatro razones.}{Tesis, caps. 1, 2 y 5.}

\NarrativeEntry{4}{Brecha}{1:00}{3:45}{
La brecha metodológica tiene tres componentes. Una comparación contra un spline propio puede ser útil para desarrollo, pero no sustituye una implementación comunitaria fuerte. Una evaluación basada solo en MAE oculta tensiones con forma temporal, fidelidad espectral y tiempo. Y una prueba con pérdidas aleatorias leves no cubre fallas localizadas. La tesis responde combinando un comparador fuerte, MNE-Python; un portafolio de métricas; y varios regímenes de pérdida. Además, separa la exploración de la evaluación final. Esa separación es crucial: la amplitud de dieciocho métodos, nueve grafos y ciento veinte escenarios orienta decisiones, pero las cifras centrales provienen únicamente del protocolo congelado TRSS fijo frente a MNE-Python.
}{Enfatizar “protocolo congelado”.}{Tesis, Introducción y Metodología.}

\NarrativeEntry{5}{Preguntas y objetivo}{1:05}{4:50}{
La primera pregunta busca las condiciones en que la regularización espacio--temporal mejora la reconstrucción frente a métodos espaciales de referencia. La segunda pregunta examina el compromiso práctico entre precisión, fidelidad espectral y tiempo de cómputo al elegir entre TRSS y MNE-Python. El objetivo general es comparar esas dimensiones bajo un marco riguroso y reproducible. No planteo una superioridad universal ni una validación clínica. La tesis busca una respuesta condicional y trazable: identificar dónde aparece la ventaja, cuantificar sus contrapartidas y explicitar qué evidencia sería necesaria para extender el alcance.
}{Leer P1 y P2 casi literalmente; luego resumir el objetivo con voz propia.}{Tesis, Introducción, preguntas P1 y P2.}

\NarrativeEntry{6}{EEG como señal sobre un grafo}{1:15}{6:05}{
El procesamiento de señales en grafos representa cada electrodo como un nodo y la relación entre electrodos como aristas ponderadas. Una muestra EEG es entonces una señal definida sobre esos nodos. El Laplaciano del grafo cuantifica variación espacial: si dos nodos conectados toman valores muy distintos, la energía de suavidad aumenta. Cuando falta un canal, este formalismo permite reconstruirlo buscando una señal compatible con los canales observados y, a la vez, regular respecto de la topología. La ventaja conceptual es que la noción de vecindad queda explícita y puede inspeccionarse. En la comparación final se privilegian grafos geométricos deterministas, construidos desde coordenadas de electrodos, para no usar la señal evaluada en la definición de sus propias relaciones.
}{No explicar toda la teoría espectral; señalar nodos, aristas y Laplaciano.}{Tesis, cap. 2 y figura de conceptos GSP.}

\NarrativeEntry{7}{MNE frente a TRSS}{1:25}{7:30}{
Los dos métodos finales usan información diferente. MNE-Python interpola mediante splines esféricos y añade decisiones de ingeniería como normalización, restricciones, pseudoinversa y manejo robusto de casos numéricos. Por eso es una referencia más exigente que un spline académico simplificado. TRSS resuelve un problema conjunto sobre canales y tiempo. El primer término obliga a respetar las observaciones; el segundo, controlado por alfa, penaliza variación espacial sobre el grafo; y el tercero, controlado por beta, penaliza cambios abruptos entre muestras. Si beta es cero, queda una regularización espacial independiente por instante. La hipótesis operativa no es que temporalidad siempre sea mejor, sino que puede aportar cuando la geometría observada deja poca información local. La misma continuidad puede transformarse en desventaja si suaviza de forma inapropiada o si su costo domina.
}{Señalar los tres términos del diagrama, no desarrollar la solución algebraica.}{Tesis, cap. 2, formulación TRSS y mecanismo de MNE.}

\NarrativeEntry{8}{Decisiones de alcance}{1:00}{8:30}{
Antes de mostrar experimentos, explicito cuatro decisiones. MNE-Python es la línea base principal para evitar comparadores debilitados. Los grafos finales son geométricos y reproducibles; los aprendidos quedan como trabajo futuro. No se incluyeron modelos entrenables porque no había un protocolo externo de entrenamiento que permitiera compararlos sin abrir otro problema experimental. Finalmente, la lectura principal usa TRSS fijo, con parámetros congelados antes de la evaluación. Estas decisiones restringen el alcance, pero también protegen la interpretación. No significan que las alternativas excluidas sean incapaces; significan que esta tesis no aporta evidencia suficiente para compararlas de manera justa.
}{Mantener tono defensivo pero no apologético.}{Tesis, Metodología y Discusión.}

\NarrativeEntry{9}{Estrategia experimental}{1:25}{9:55}{
El flujo experimental tiene varias capas. Primero se estandarizan y segmentan tres bases públicas. Después se generan máscaras de pérdida con semillas controladas. Luego se construyen grafos, se ejecutan métodos y se calculan métricas. La fase exploratoria recorre dieciocho implementaciones y nueve grafos; la fase intermedia reduce candidatos por estabilidad, costo y coherencia; después se optimizan parámetros y se congelan variantes. Solo entonces se ejecuta la comparación final pareada. Esta secuencia evita confundir búsqueda con confirmación. También permite que cada resultado conserve un identificador, la máscara utilizada, parámetros, tiempos y señales reconstruidas. La trazabilidad no es un apéndice administrativo: es lo que permite regresar desde una cifra de la defensa hasta la configuración que la originó.
}{Recorrer visualmente el flujo y detenerse en el congelamiento.}{Tesis, figura de flujo metodológico.}

\NarrativeEntry{10}{Datos}{1:05}{11:00}{
El estudio usa tres conjuntos públicos con características complementarias. PhysioNet aporta 64 canales, 160 hertz y 109 sujetos en reposo. BCI Competition IV 2a aporta 22 canales, 250 hertz y nueve sujetos en imaginación motora. MNE Sample aporta 60 canales, 600 hertz y un sujeto en un paradigma auditivo. Esta diversidad prueba montajes, frecuencias y paradigmas diferentes. Sin embargo, no debe interpretarse como una muestra representativa de toda población EEG. Especialmente MNE Sample contribuye diversidad técnica, no evidencia poblacional. La conclusión válida es que el comportamiento se observó bajo varios entornos públicos controlados; extenderlo a otros equipos, montajes o poblaciones exige validación adicional.
}{Contrastar explícitamente diversidad técnica con generalización poblacional.}{Tesis, Metodología, tabla de conjuntos de datos.}

\NarrativeEntry{11}{Cómo se simula la pérdida}{1:05}{12:05}{
Se estudian dos patrones. En el aleatorio, los canales faltantes se distribuyen; en el contiguo, se elige un canal semilla y se agregan vecinos cercanos, aproximando una falla localizada. La fase amplia cruza tres datasets, dos modos, cuatro severidades y cinco semillas: ciento veinte escenarios. Dentro de cada escenario, todos los métodos reciben exactamente la misma señal y la misma máscara. Esto permite comparaciones pareadas y evita que un método se beneficie de una pérdida más sencilla. La ventaja del ocultamiento artificial es conocer la verdad; su límite es que un canal real puede tener ruido residual, deriva o artefactos, no solo ausencia completa.
}{Explicar primero el dibujo mental de las máscaras y luego la aritmética 3 por 2 por 4 por 5.}{Tesis, Metodología, generación de máscaras.}

\NarrativeEntry{12}{Evaluación multi-métrica}{1:05}{13:10}{
La evaluación separa dimensiones. MAE, RMSE y NRMSE miden error de amplitud; SNR la relación entre señal y error; DTW, correlación y R cuadrado capturan aspectos de forma temporal; LSD, coherencia y PSD examinan el dominio espectral; y el tiempo caracteriza costo. Cada métrica tiene una dirección favorable definida antes del análisis. No construyo un puntaje único, porque eso obligaría a imponer pesos subjetivos entre objetivos distintos. En su lugar, presento un portafolio y pregunto si la misma conclusión se sostiene en cada eje. Esta elección permite reconocer que un método puede mejorar la amplitud y, simultáneamente, perder en fidelidad espectral o velocidad.
}{Nombrar familias, no recitar fórmulas.}{Tesis, Metodología, sección Métricas.}

\NarrativeEntry{13}{Arquitectura y trazabilidad}{1:00}{14:10}{
La arquitectura implementa esa lógica con configuración inmutable, carga de datos, construcción del grafo, ejecución del interpolador, cálculo de métricas y persistencia de artefactos. Cada caso conserva dataset, patrón, severidad, semilla, máscara, parámetros, tiempos y reconstrucción. Un identificador SHA-256 permite detectar cambios y vincular filas, figuras y configuraciones. También se validan dimensionalidad, canales observados y faltantes, coherencia temporal y correspondencia entre máscara y semilla. El objetivo es que una cifra no dependa de memoria manual ni de una captura aislada: debe poder localizarse, regenerarse y auditarse.
}{Señalar el recorrido entrada--resultado; no entrar a detalles de código.}{Tesis, Metodología y figura de arquitectura.}

\NarrativeEntry{14}{Reducción experimental}{1:10}{15:20}{
La amplitud inicial no equivale a la comparación final. La exploración consideró dieciocho implementaciones agrupadas en métodos geométricos, estadísticos, GSP y espacio--temporales, además de nueve grafos. Luego se redujeron candidatos por precisión, estabilidad numérica, costo y reproducibilidad. Algunas opciones aprendidas se descartaron por costo y riesgo de circularidad; los métodos geométricos simples sirvieron como referencia exploratoria; TRSS avanzó por su potencial espacio--temporal. La fase de optimización ajustó familias, pero la evaluación final congeló TRSS fijo y MNE-Python. La variante calibrada y el oráculo se conservan solo para análisis auxiliar. Por eso no mezclo cifras exploratorias con las que responden P1 y P2.
}{Usar la caída de 18 métodos a 2 como argumento de rigor, no como ranking total.}{Tesis, tablas de fases y decisiones de reducción.}

\NarrativeEntry{15}{Protocolo final}{1:25}{16:45}{
El protocolo final contiene cien estratos definidos por régimen, patrón y severidad. En cada estrato hay tres semillas, para un total de trescientas diferencias pareadas TRSS--MNE. Dentro de cada par se comparten señal, máscara y preprocesamiento; la diferencia se calcula en el mismo caso. La unidad observacional es esa diferencia por estrato y semilla, con $n=300$; la unidad de agrupación es el estrato, con $n=100$. Para resumir estabilidad se usa un bootstrap jerárquico descriptivo de mil remuestreos que respeta ese agrupamiento. Las semillas no son sujetos nuevos. Los intervalos describen la variabilidad del protocolo; no son pruebas de hipótesis ni justifican inferencia poblacional externa. Esta precisión evita inflar el tamaño efectivo de la evidencia.
}{Enfatizar 100, 3 y 300; luego aclarar qué no significa el intervalo.}{Tesis, Protocolo de Validación Comparativa.}

\NarrativeEntry{16}{Resultado principal en MAE}{1:30}{18:15}{
El resultado central es una mejora mediana de 12,4 por ciento en MAE para TRSS fijo frente a MNE-Python. El intervalo descriptivo al 95 por ciento va de 7,8 a 17,4 por ciento y TRSS gana en 72 por ciento de los trescientos casos. Esta es la evidencia más directa a favor de la regularización temporal bajo el protocolo congelado. Pero la cifra no significa que cada escenario mejore ni que la diferencia se transfiera automáticamente a otro montaje. Resume el conjunto de casos diseñados aquí. La tasa de victoria complementa la mediana: muestra que la tendencia no depende de un pequeño número de mejoras extremas, aunque todavía existe una proporción relevante de empates o victorias de MNE.
}{Pausa después de 12,4 por ciento. Señalar intervalo y tasa de victoria.}{Tesis, tabla principal robusta y figura de intervalos.}

\NarrativeEntry{17}{Portafolio de métricas}{1:30}{19:45}{
Cuando ampliamos la mirada, TRSS conserva ventajas agregadas en NRMSE, DTW, SNR, correlación y R cuadrado. Por ejemplo, la mejora mediana es 13,9 por ciento en NRMSE y 9,8 por ciento en DTW. Sin embargo, LSD muestra una mejora mediana de menos 2,3 por ciento y una tasa de victoria TRSS de 40,7 por ciento: en esa métrica global, MNE es mejor. El tiempo también favorece claramente a MNE. Esta tensión es parte del resultado, no una anomalía que deba ocultarse. La respuesta correcta a P2 no es declarar un ganador, sino reconocer una frontera: TRSS mejora varios aspectos de amplitud y forma temporal, mientras MNE preserva mejor el LSD agregado y ejecuta mucho más rápido.
}{Mostrar primero el grupo favorable y después detenerse en LSD y tiempo.}{Tesis, tabla y figura del portafolio de métricas.}

\NarrativeEntry{18}{Heterogeneidad por escenario}{1:30}{21:15}{
La ventaja no se distribuye uniformemente. En pérdidas cercanas o agrupadas con 30 y 40 por ciento de canales faltantes, la mejora mediana en MAE supera 40 por ciento y la tasa de victoria alcanza 80 u 86,7 por ciento. En pérdidas periféricas severas también aparece una ventaja importante. En cambio, con dos canales cercanos faltantes, la mediana puede ser cero y la tasa de victoria queda próxima al empate. La interpretación mecánica es coherente: cuando desaparecen vecinos informativos, la continuidad temporal aporta una restricción adicional; cuando la geometría espacial conserva suficiente información, el costo y el suavizado extra pueden no aportar. Estos estratos ilustran condiciones frontera, pero no estiman con qué frecuencia aparecerán en una práctica externa.
}{Comparar un extremo favorable y el caso de empate; evitar leer toda la matriz.}{Tesis, tabla de escenarios seleccionados y mapa de calor.}

\NarrativeEntry{19}{Forma de señal}{1:10}{22:25}{
Las series temporales representativas ayudan a entender el mecanismo, pero no reemplazan el agregado. Se observan casos donde TRSS sigue mejor la forma original, otros donde MNE es suficiente y algunos donde ambos fallan de manera distinta. Las curvas permiten detectar suavizado, desfase o desviaciones locales que un número resume parcialmente. Por eso se usan como ejemplos auditables y no como prueba de superioridad. La conclusión cuantitativa sigue descansando en los trescientos pares y en el análisis por estratos. El valor de esta lámina es conectar la estadística con una señal concreta y recordar qué significa, en términos físicos, reducir el error.
}{Señalar original y reconstrucciones; no inferir desde un único canal.}{Tesis, Resultados, series temporales representativas.}

\NarrativeEntry{20}{Tensión espectral}{1:10}{23:35}{
El análisis de densidad espectral muestra por qué era necesario un portafolio. Un método puede aproximar mejor la señal muestra a muestra y, aun así, alterar la distribución de energía entre frecuencias. El LSD agregado favorece a MNE: TRSS tiene una mejora mediana negativa de 2,3 por ciento y gana solo en 40,7 por ciento de los casos. Esto no invalida el resultado en MAE; indica que los objetivos no son intercambiables. Si la prioridad de una aplicación es preservar el espectro global, MNE tiene una ventaja bajo este benchmark. Si importa la reconstrucción temporal en una pérdida localizada y el costo es aceptable, TRSS puede ser preferible.
}{Conectar explícitamente PSD cualitativa con LSD cuantitativo.}{Tesis, Resultados, PSD y portafolio de métricas.}

\NarrativeEntry{21}{Costo computacional}{1:05}{24:40}{
La diferencia de costo es clara. MNE-Python presenta una mediana de 0,0090 segundos por caso; TRSS fijo, 0,1214 segundos. En esta implementación y hardware, TRSS es aproximadamente un orden de magnitud más lento. Los tiempos no son constantes universales: dependen del solver, tamaño, hardware e ingeniería. Aun así, la dirección de la diferencia es suficientemente estable para formar parte de la decisión. La variante calibrada no elimina el costo de búsqueda de hiperparámetros y no es la lectura principal. Por tanto, un sistema de alto volumen o con restricción de latencia debe privilegiar MNE, salvo que la dificultad del patrón de pérdida justifique el costo adicional.
}{Dar cifras con cuatro decimales y aclarar dependencia de implementación.}{Tesis, tabla de tiempo y complejidad.}

\NarrativeEntry{22}{Frontera práctica}{1:05}{25:45}{
La evidencia se traduce en una regla condicional. MNE es la opción por defecto cuando faltan pocos canales, la geometría conserva vecinos útiles, se prioriza LSD global o existe una restricción fuerte de tiempo. TRSS merece consideración cuando la pérdida es espacialmente difícil, faltan vecinos cercanos o periféricos, el análisis temporal es importante y el procesamiento fuera de línea admite mayor costo. Esta regla no automatiza una decisión clínica ni reemplaza validación posterior. Es un mapa de diseño basado en el benchmark: permite elegir de manera explícita qué objetivo se prioriza, en vez de asumir que una sola métrica decide por todos.
}{Recorrer los dos lados del mapa y cerrar en “regla condicional”.}{Tesis, Discusión y mapa de decisión.}

\NarrativeEntry{23}{Aportes}{1:00}{26:45}{
El trabajo aporta cuatro elementos. Primero, una caracterización condicional de cuándo la regularización temporal agrega valor. Segundo, una metodología que separa exploración, reducción, congelamiento y evaluación final. Tercero, una comparación pareada y multi-métrica contra una referencia comunitaria fuerte. Cuarto, una cadena reproducible desde la configuración y la máscara hasta la tabla, la figura y la afirmación. La novedad no descansa en inventar el procesamiento de grafos ni los splines, sino en integrar TRSS, un benchmark exigente y una lectura de frontera respaldada por trazabilidad.
}{Enumerar con la mano; no sobredimensionar la novedad algorítmica.}{Tesis, Conclusiones, sección Aportes.}

\NarrativeEntry{24}{Limitaciones}{1:15}{28:00}{
Las limitaciones definen qué no concluyo. El ocultamiento artificial entrega verdad de terreno, pero no reproduce ruido residual ni deterioro real. Las medidas repetidas permiten comparación pareada, pero no equivalen a sujetos independientes ni sostienen inferencia poblacional. Las métricas son proxies de reconstrucción; falta evaluar impacto en clasificación, conectividad o lectura experta. Los casos visuales no son una evaluación ciega. Los tiempos pertenecen a una implementación y hardware reportados. Tampoco se compararon modelos entrenables bajo un protocolo externo. El siguiente paso debe incluir daño real, tareas posteriores, evaluación ciega, otros montajes y una versión acelerada de TRSS. Presentar estos límites no debilita la tesis: evita que una evidencia controlada se transforme en una promesa que no fue ensayada.
}{Reducir velocidad y mirar al jurado; esta lámina prepara preguntas difíciles.}{Tesis, Discusión y Conclusiones.}

\NarrativeEntry{25}{Conclusión}{1:00}{29:00}{
Concluyo con dos respuestas. Para P1, TRSS fijo mejora la reconstrucción en el agregado y aporta especialmente cuando la pérdida espacial elimina vecinos informativos; no reemplaza universalmente a MNE. Para P2, la ganancia en amplitud y forma temporal se intercambia por mayor costo y una desventaja en LSD global. La decisión depende del régimen y del objetivo posterior. El aporte central es haber identificado esa frontera con un protocolo pareado, congelado y trazable, mostrando tanto la ventaja como sus límites. Muchas gracias; quedo atento a sus preguntas.
}{Pausa entre P1 y P2. Terminar mirando al jurado, no a la pantalla.}{Tesis, Conclusiones y respuestas a P1/P2.}

\clearpage
\section*{Banco de preguntas del jurado}
\addcontentsline{toc}{section}{Banco de preguntas del jurado}

\subsection*{A. Estadística y diseño experimental}

\QA{1. ¿Por qué no realizó pruebas de hipótesis tradicionales?}
{Porque la estructura repetida no equivale a una muestra aleatoria de sujetos independientes. Reporto diferencias pareadas, medianas, tasas de victoria e intervalos descriptivos.}
{Los 300 casos provienen de 100 estratos y tres semillas por estrato. Tratar cada máscara como una observación poblacional independiente inflaría la evidencia. El bootstrap jerárquico conserva el agrupamiento y cuantifica estabilidad dentro del protocolo, pero no transforma el benchmark en una muestra poblacional. Por eso las conclusiones son descriptivas y condicionales.}
{Metodología, Protocolo de Validación Comparativa.}

\QA{2. ¿Cuál es exactamente la unidad estadística?}
{El caso pareado TRSS--MNE dentro de un estrato; no la semilla como sujeto independiente.}
{Cada par comparte señal, máscara, semilla y preprocesamiento. La diferencia se calcula dentro del caso y los casos están agrupados en estratos de régimen, patrón y severidad. Las semillas generan repeticiones controladas, no nuevos participantes. Esa estructura se respeta en el resumen jerárquico.}
{Metodología, unidad pareada y bootstrap jerárquico.}

\QA{3. ¿Por qué usar la mediana?}
{Porque las mejoras porcentuales y los tiempos son asimétricos y sensibles a casos extremos.}
{La mediana ofrece un resumen robusto del comportamiento típico sin quedar dominada por denominadores pequeños o escenarios de error extremo. No se usa sola: se acompaña de intervalo descriptivo, tasa de victoria y análisis por estrato para mostrar dispersión y heterogeneidad.}
{Tablas robustas de Resultados.}

\QA{4. ¿Qué significa el intervalo de 7,8 a 17,4 por ciento en MAE?}
{Describe la estabilidad de la mejora mediana bajo el remuestreo jerárquico del protocolo.}
{No es un intervalo poblacional ni una prueba de significancia. Resume cómo varía la mediana al remuestrear estratos y casos respetando su estructura. Su utilidad es mostrar que la dirección positiva en MAE no depende de una única partición de los 300 pares.}
{Tabla principal robusta.}

\QA{5. ¿Por qué mil remuestreos?}
{Es un compromiso reproducible entre estabilidad descriptiva y costo computacional.}
{Mil réplicas permiten estimar cuantiles con resolución suficiente para esta lectura descriptiva. Aumentar el número puede refinar ligeramente los extremos, pero no corrige sesgos de diseño ni amplía la validez externa. La prioridad fue respetar la jerarquía y documentar la semilla.}
{Metodología, bootstrap jerárquico.}

\QA{6. ¿Cómo evitó comparar métodos con máscaras distintas?}
{La máscara y la semilla se generan una vez por caso y se reutilizan para todos los métodos.}
{La configuración persistida identifica dataset, segmento, patrón, severidad y semilla. Los validadores comprueban canales observados y faltantes, dimensionalidad y correspondencia de la máscara. Así, la diferencia TRSS--MNE se atribuye al método y no a una realización distinta de la pérdida.}
{Arquitectura y controles de calidad.}

\QA{7. ¿Los 300 pares son 300 sujetos?}
{No. Son casos experimentales repetidos y agrupados; no deben interpretarse como 300 sujetos.}
{Los casos cruzan regímenes, patrones, severidades y semillas. Algunas bases aportan múltiples sujetos y otras, como MNE Sample, un solo sujeto. Por eso la tesis evita lenguaje de inferencia poblacional y reporta estabilidad del benchmark, no prevalencia clínica.}
{Metodología y Limitaciones.}

\QA{8. ¿Cómo trató comparaciones múltiples?}
{No se basó la conclusión en una batería de p-valores; se separaron dimensiones y se declaró una lectura descriptiva.}
{Cada métrica responde una pregunta distinta y no se colapsa en un ranking universal. La tesis presenta direcciones, intervalos descriptivos y tasas de victoria, y reconoce tensiones como LSD y tiempo. Si se diseñara un estudio confirmatorio con hipótesis múltiples, habría que preespecificar una métrica primaria y controlar multiplicidad.}
{Portafolio de métricas y Discusión.}

\subsection*{B. MNE-Python, TRSS y grafos}

\QA{9. ¿Por qué MNE-Python es la línea base principal?}
{Porque representa una implementación comunitaria madura y evita una comparación contra un spline debilitado.}
{MNE incorpora normalización, restricciones, pseudoinversa y manejo robusto de casos numéricos. Un spline propio es útil para aislar componentes, pero no equivale a esa ingeniería. Usar MNE hace más exigente y más útil la comparación principal.}
{Marco Teórico, mecanismo de MNE.}

\QA{10. ¿Por qué no comparar solo con su spline propio?}
{Porque habría reducido la fuerza externa del benchmark.}
{Una implementación propia puede diferir por detalles numéricos y de preprocesamiento. La tesis conserva esas variantes en exploración, pero las cifras centrales responden frente a MNE-Python. Así se evita atribuir a TRSS una ventaja causada por una referencia incompleta.}
{Introducción y decisiones de reducción.}

\QA{11. ¿Cuál es el papel de alfa y beta en TRSS?}
{Alfa controla suavidad espacial; beta controla continuidad temporal.}
{El objetivo combina fidelidad en canales observados, energía espacial sobre el Laplaciano y diferencias temporales finitas. Con beta igual a cero se obtiene una reconstrucción espacial independiente por instante; con alfa igual a cero se pierde la guía del grafo. La variante final fija ambos antes de evaluar.}
{Marco Teórico, ecuación TRSS.}

\QA{12. ¿TRSS puede sobre-suavizar la señal?}
{Sí; esa es una de sus contrapartidas y ayuda a explicar la tensión espectral.}
{La penalización temporal restringe cambios muestra a muestra. Si beta es demasiado alto o la señal contiene transitorios relevantes, puede disminuir error agregado y alterar detalles. Por eso se evalúan DTW, correlación, PSD y LSD, y no solo MAE.}
{Marco Teórico y Resultados espectrales.}

\QA{13. ¿Por qué grafos geométricos y no funcionales?}
{Por determinismo, reproducibilidad y para evitar usar la señal evaluada al construir el comparador.}
{Los grafos aprendidos pueden capturar relaciones funcionales, pero requieren un protocolo separado: datos de entrenamiento, validación externa y control de circularidad. En esta tesis se exploraron, pero la fase final privilegió grafos construidos desde coordenadas conocidas.}
{Metodología, construcción y selección de grafos.}

\QA{14. ¿Cómo eligió el grafo final?}
{Mediante exploración y optimización separadas de la evaluación final, seguida de congelamiento.}
{Se compararon nueve familias y se redujeron por precisión, estabilidad, costo y reproducibilidad. La configuración final no se reelige por cada caso de evaluación. Esa regla es importante para que la mejora no sea un oráculo posterior.}
{Fases experimentales y decisiones de reducción.}

\QA{15. ¿Qué ocurre si beta es cero?}
{TRSS se reduce a una regularización espacial tipo Tikhonov aplicada por instante.}
{Ese caso permite interpretar el aporte marginal de temporalidad. La tesis usa la descomposición para entender el método, pero la cifra principal corresponde a la configuración fija completa, no a seleccionar beta después de observar cada resultado.}
{Marco Teórico y estudio de descomposición.}

\QA{16. ¿Por qué no usar un grafo dinámico?}
{Porque habría agregado aprendizaje y costo sin un protocolo externo suficiente para evaluarlo limpiamente.}
{Un grafo dinámico podría adaptarse a estados temporales, pero también puede usar información del mismo segmento que luego se evalúa. Para sostenerlo harían falta particiones estrictas, estabilidad temporal y análisis de costo. Se propone como trabajo futuro, no como ausencia de valor.}
{Discusión y trabajo futuro.}

\QA{17. ¿La ventaja se debe solo a más hiperparámetros?}
{No se puede reducir a ese hecho; la lectura principal usa una configuración fija y la ventaja aparece por régimen.}
{TRSS tiene alfa y beta, pero no se seleccionan por caso en el protocolo central. Además, la ventaja crece cuando faltan vecinos y se reduce en casos leves, un patrón coherente con el uso de información temporal. Aun así, un estudio de ablación más amplio podría aislar mejor cada término.}
{Resultados por escenario y Limitaciones.}

\subsection*{C. Métricas y resultados}

\QA{18. ¿Por qué MAE es la métrica principal?}
{Porque es interpretable en amplitud y robusta, pero no se usa como único criterio.}
{MAE mide desviación absoluta en los canales ocultos sin elevar errores extremos al cuadrado. La tesis la usa como eje central y la acompaña con NRMSE, DTW, SNR, correlación, LSD, coherencia y tiempo. La conclusión se formula como portafolio, no como monopolio de MAE.}
{Metodología y tabla robusta.}

\QA{19. ¿Qué significa una mejora de 12,4 por ciento?}
{Que la reducción relativa mediana del MAE de TRSS respecto de MNE es 12,4 por ciento en el protocolo final.}
{La mejora se calcula por caso con una orientación común donde valores positivos favorecen a TRSS y luego se resume de forma robusta. No significa que todos los casos mejoren: TRSS gana en 72 por ciento, por lo que existe heterogeneidad relevante.}
{Tabla principal robusta.}

\QA{20. ¿Por qué MNE gana en LSD?}
{Porque menor error temporal no garantiza preservar mejor la distribución espectral global.}
{La regularización temporal puede suavizar o redistribuir energía. MNE, basado en interpolación espacial por instante, muestra una mediana favorable en LSD y mayor tasa de victoria. El resultado exige elegir según el objetivo posterior y justifica no afirmar superioridad universal de TRSS.}
{Portafolio de métricas y PSD.}

\QA{21. ¿Puede un método ganar en MAE y perder en SNR?}
{Sí, porque las métricas normalizan y ponderan el error de manera distinta.}
{MAE resume magnitud absoluta, mientras SNR relaciona energía de señal y error. En este estudio ambos favorecen agregadamente a TRSS, pero no deben suponerse equivalentes caso a caso. Esa diferencia motiva reportar el portafolio completo.}
{Definición de métricas.}

\QA{22. ¿Por qué incluir DTW?}
{Para evaluar similitud de forma temporal permitiendo pequeñas desalineaciones.}
{DTW complementa MAE porque una señal puede conservar una trayectoria parecida aun con desplazamientos locales. Tampoco es suficiente por sí sola: puede tolerar deformaciones temporales, por lo que se interpreta junto a correlación, error y espectro.}
{Metodología, métricas temporales.}

\QA{23. ¿Los casos visuales prueban la ventaja?}
{No. Son ejemplos interpretables; la conclusión descansa en los 300 pares.}
{Las series y PSD permiten detectar forma, suavizado y distribución espectral. Fueron elegidas para ilustrar comportamientos favorables, desfavorables y de empate. No reemplazan la estadística agregada ni una evaluación ciega.}
{Resultados y Limitaciones.}

\QA{24. ¿Cuál es el escenario más favorable a TRSS?}
{Las pérdidas cercanas o agrupadas severas, especialmente 30 y 40 por ciento.}
{En los escenarios seleccionados, la mejora mediana en MAE supera 40 por ciento y la tasa de victoria llega a 80 u 86,7 por ciento. La explicación propuesta es la falta de vecinos espaciales informativos, donde la continuidad temporal aporta una restricción adicional.}
{Tabla de escenarios seleccionados.}

\QA{25. ¿Dónde no aporta TRSS?}
{En pérdidas leves con suficiente información espacial puede empatar o perder.}
{Con dos canales cercanos faltantes, la tabla seleccionada muestra mejora mediana cercana a cero y una tasa de victoria inferior a la mitad. Además, MNE domina en tiempo y LSD agregado. Por tanto, TRSS no es una elección automática.}
{Escenarios seleccionados y portafolio.}

\QA{26. ¿Por qué no combinar todas las métricas en un puntaje?}
{Porque requeriría pesos de utilidad que el experimento no justifica.}
{Un puntaje único ocultaría que amplitud, espectro y costo responden preguntas diferentes. En una aplicación concreta se podrían definir pesos a priori, pero para una tesis comparativa es más transparente mostrar la frontera y dejar explícita la prioridad.}
{Discusión multi-métrica.}

\subsection*{D. Validez, alcance y reproducibilidad}

\QA{27. ¿Qué tan reproducible es el estudio?}
{Cada caso conserva configuración, semilla, máscara, métricas, tiempos, reconstrucción e identificador.}
{La arquitectura permite rastrear una figura o tabla hasta resultados derivados y configuraciones. La reproducibilidad todavía depende de versiones de software y hardware, por lo que estos también deben registrarse. La defensa conserva una matriz claim--evidence y checksums de fuentes.}
{Metodología, trazabilidad; artefactos de defensa.}

\QA{28. ¿Cómo controló fuga de información?}
{La verdad del canal oculto se usa para evaluar, no para construir la máscara ni ajustar por caso el método final.}
{Los grafos finales son geométricos y la configuración principal está congelada. Las mismas observaciones disponibles alimentan a ambos métodos. Un oráculo de parámetros se reporta solo como cota auxiliar y no como resultado desplegable.}
{Metodología y decisiones de alcance.}

\QA{29. ¿Puede afirmar utilidad clínica?}
{No. El estudio evalúa reconstrucción controlada, no desenlaces ni decisiones clínicas.}
{Faltan canales realmente dañados, evaluación ciega por especialistas, impacto en tareas posteriores, diversidad de equipos y un protocolo prospectivo. La tesis ofrece evidencia metodológica que puede motivar esos estudios, pero no valida uso clínico rutinario.}
{Discusión, limitaciones.}

\QA{30. ¿Por qué tres datasets son suficientes?}
{Son suficientes para el alcance controlado de esta tesis, no para generalización universal.}
{Aportan diversidad de montajes, tasas y paradigmas, pero sus tamaños y estructuras son heterogéneos. La conclusión se limita al benchmark. Un estudio externo debería añadir más bases, equipos, poblaciones y protocolos de daño real.}
{Metodología y Validez externa.}

\QA{31. ¿Cuál es la principal amenaza a la validez externa?}
{Que la pérdida artificial y los datasets públicos no representan toda falla real ni toda población.}
{Un electrodo real puede presentar ruido, saturación, deriva o contacto intermitente. Además, la composición de las bases no permite una muestra poblacional equilibrada. La siguiente validación debe incluir daño observado y tareas funcionales.}
{Discusión y Conclusiones.}

\QA{32. ¿Qué haría primero como trabajo futuro?}
{Validaría daño real y tareas posteriores, en paralelo con acelerar TRSS.}
{La prioridad científica es comprobar si la mejora de reconstrucción se traduce en mejor clasificación, conectividad o lectura experta bajo canales realmente dañados. La prioridad de ingeniería es reducir latencia y memoria. Solo después tendría sentido automatizar la regla de decisión.}
{Conclusiones, trabajo futuro.}

\QA{33. ¿Por qué no incluyó redes neuronales?}
{Porque requerían un protocolo de entrenamiento y generalización externo que excedía el diseño congelado.}
{Compararlas justamente exige separar sujetos o bases para entrenamiento, validación y prueba, controlar arquitectura y presupuesto, y evitar que vean patrones del conjunto final. La tesis no afirma que sean inferiores; delimita una comparación entre métodos no entrenables y reproducibles.}
{Metodología y Limitaciones.}

\QA{34. ¿Los tiempos pueden generalizarse a otro computador?}
{No como valores absolutos; sí como evidencia relativa de esta implementación reportada.}
{Los tiempos dependen de CPU, memoria, solver, tamaño de segmento y versiones. La diferencia observada, 0,0090 frente a 0,1214 segundos por caso, establece una contrapartida práctica en el entorno probado. Otro sistema debe volver a medirla.}
{Tabla de tiempo y complejidad.}

\QA{35. ¿Cuál es la contribución novedosa si TRSS y MNE ya existían?}
{La contribución es la caracterización experimental condicional, no reclamar un algoritmo completamente nuevo.}
{El trabajo integra un benchmark fuerte, pérdida estructurada, evaluación pareada multi-métrica, separación entre exploración y congelamiento, y trazabilidad. Su resultado novedoso es una frontera defendible de decisión: cuándo la temporalidad agrega valor y qué costo o degradación acompaña esa ventaja.}
{Conclusiones, Aportes.}

\QA{36. ¿Por qué la exploración tiene 120 escenarios y el protocolo final 100 estratos?}
{Porque pertenecen a fases distintas y no comparten denominador: los 120 escenarios sirven para explorar; los 100 estratos definen la evaluación final congelada.}
{La fase amplia cruza tres datasets, dos modos, cuatro severidades y cinco semillas. Esa amplitud orienta selección y congelamiento, pero no sostiene la cifra central. La comparación final usa un inventario independiente de 100 estratos y tres realizaciones por estrato, que generan 300 diferencias pareadas. Mezclar ambos conteos convertiría decisiones exploratorias en evidencia final.}
{Metodología; fases experimentales; Protocolo de Validación Comparativa.}

\QA{37. ¿El tiempo reportado incluye la búsqueda de hiperparámetros?}
{No. Los tiempos centrales corresponden a la ejecución por caso de las variantes ya fijadas; la calibración es un costo adicional y se mantiene fuera de la comparación principal.}
{MNE registra una mediana de 0,0090 segundos por caso y TRSS fijo de 0,1214 segundos en el entorno evaluado. Esas cifras comparan inferencia con configuraciones congeladas. Una variante calibrada debe sumar el costo de búsqueda, por lo que no puede presentarse como igual de desplegable ni usarse para rebajar artificialmente el costo de TRSS fijo.}
{Resultados, tabla de tiempo y complejidad; decisiones de alcance.}

\QA{38. ¿El mapa de decisión ya puede automatizar la elección del método?}
{No. Es una síntesis descriptiva de los regímenes evaluados, no una regla predictiva validada externamente.}
{El mapa traduce el benchmark a criterios transparentes: dificultad espacial, prioridad espectral, costo y procesamiento fuera de línea. Para automatizarlo habría que predefinir umbrales, evaluarlos en datos externos con daño real y medir el efecto en la tarea posterior. Hasta entonces orienta una decisión humana, pero no reemplaza validación específica.}
{Discusión, mapa de decisión y Limitaciones.}

\end{document}
